% ============================================================
% CHƯƠNG 5: KẾT LUẬN VÀ HƯỚNG PHÁT TRIỂN
% ============================================================
\chapter{Kết luận và hướng phát triển}
\label{chap:conclusion}

\section{Kết luận}
\label{sec:conclusion}

Đồ án đã hoàn thành đầy đủ và toàn diện các mục tiêu nghiên cứu đề ra, đóng góp các kết quả cụ thể trên ba phương diện chính bao gồm phân tích dữ liệu, cải tiến mô hình và thực nghiệm đánh giá hiệu năng.

Về phương diện phân tích dữ liệu, đồ án đã tiến hành phân tích chuyên sâu bộ dữ liệu chuẩn IEMOCAP với quy mô 2.280 câu nói thuộc 4 lớp cảm xúc chính. Nghiên cứu đã chỉ ra sự mất cân bằng dữ liệu cực đoan giữa các lớp, với lớp Neutral chiếm đa số đạt 48,2\% và lớp Happiness chỉ đạt tỷ lệ rất thấp là 12,5\%. Đồng thời, quy trình tiền xử lý tín hiệu hoàn chỉnh cũng được xây dựng thành công, chuyển đổi từ tín hiệu âm thanh thô ban đầu sang phổ đồ Log-Spectrogram, định dạng ảnh Pseudo-RGB và đưa về tensor đầu vào chuẩn cho mạng CNN. Phân tích cũng đã làm rõ nguyên nhân nhầm lẫn phổ biến giữa các cảm xúc Neutral và Sadness dưới góc độ vật lý của tín hiệu âm thanh.

Về phương diện cải tiến mô hình, nghiên cứu đã triển khai thành công 6 kiến trúc mạng CNN tiên tiến bao gồm AlexNet, AlexNet GAP, ResNet-18, DenseNet-121, EfficientNet-B0 và MobileNet-V2, kết hợp đồng bộ kỹ thuật Transfer Learning từ ImageNet để nâng cao khả năng học đặc trưng. Đồng thời, hàm lỗi Focal Loss với hệ số tập trung $\gamma = 2.0$ và cơ chế trọng số động cũng được tích hợp để giải quyết hiện tượng mất cân bằng lớp. Bên cạnh đó, ba biến thể tăng cường dữ liệu EMix gồm EMix-S, EMix-N và EMix-NS cũng được đưa vào nhằm tăng tính bền vững của mô hình.

Về kết quả thực nghiệm, cấu hình kết hợp tối ưu nhất bao gồm EfficientNet-B0, hàm lỗi Focal Loss và tăng cường dữ liệu EMix-S đã đạt hiệu năng vượt trội. Cụ thể, trên lượt kiểm thử đơn lẻ, mô hình đạt độ chính xác WA đạt 80,49\% và UA đạt 73,25\%. Trên quy trình đánh giá chéo 5-Fold, mô hình duy trì kết quả trung bình WA đạt $69,80\% \pm 7,29\%$ và UA đạt $62,27\% \pm 8,18\%$. Mô hình đã cải thiện mạnh mẽ độ chính xác trên lớp thiểu số nhất Happiness từ 0,00\% ở baseline lên mức 61,54\% và giảm tới 93\% lượng tham số so với AlexNet gốc. Kết quả thực nghiệm Ablation Study cũng khẳng định vai trò bổ trợ chặt chẽ của cả Focal Loss và EMix-S đối với sự nâng cao hiệu suất chung.

\subsection{Tổng kết đóng góp chính}

Bảng~\ref{tab:contribution} tổng hợp mức đóng góp riêng lẻ của từng kỹ thuật cải tiến trong hệ thống đối với hiệu năng nhận dạng cảm xúc chung.

\begin{table}[H]
\centering
\caption{Tổng kết đóng góp của từng kỹ thuật cải tiến}
\label{tab:contribution}
\begin{tabular}{llll}
\toprule
\textbf{Kỹ thuật} & \textbf{Vai trò} & \textbf{Cải thiện WA} & \textbf{Cải thiện UA} \\
\midrule
AlexNet GAP & Giảm overfitting & +27,81\% & +39,89\% \\
Focal Loss & Xử lý mất cân bằng & +$\sim$4--6\% & +$\sim$10\% \\
EMix-S & Tăng cường dữ liệu & +$\sim$3--5\% & +$\sim$5--8\% \\
EfficientNet-B0 & Kiến trúc tối ưu & +$\sim$2--4\% & +$\sim$2--5\% \\
\bottomrule
\end{tabular}
\end{table}


\section{Hạn chế}
\label{sec:limitations}

Mặc dù đạt được những kết quả khả quan, đồ án vẫn tồn tại một số hạn chế kỹ thuật nhất định. Thứ nhất, hệ thống mới chỉ được thử nghiệm trên bộ dữ liệu tiếng Anh IEMOCAP mà chưa đánh giá trên các ngôn ngữ khác, đặc biệt là tiếng Việt vốn có đặc tính thanh điệu phức tạp. Thứ hai, bài toán phân loại hiện tại mới chỉ giới hạn trên 4 lớp cảm xúc cơ bản, chưa mở rộng sang các trạng thái cảm xúc phức tạp hoặc đa dạng hơn. Thứ ba, độ lệch chuẩn trong đánh giá chéo 5-Fold vẫn còn ở mức cao (khoảng 7--8$\%$), cho thấy mô hình còn nhạy cảm với sự biến thiên phong cách biểu biểu đạt của từng cá nhân. Thứ tư, về mặt đặc trưng, đồ án mới chỉ khai thác đặc trưng Log-Spectrogram thô mà chưa kết hợp các đặc trưng âm học truyền thống giàu thông tin như MFCC, cao độ (pitch) hay năng lượng (energy). Cuối cùng, hệ thống chưa tận dụng cấu trúc thông tin đa phương thức (multimodal) bao gồm kênh văn bản hội thoại (transcriptions) và kênh video có sẵn trong bộ dữ liệu.


\section{Hướng phát triển}
\label{sec:future_work}

Từ các hạn chế nêu trên, các hướng phát triển tiếp theo của đề tài được vạch ra rõ ràng. Định hướng thứ nhất là mở rộng phạm vi kiểm thử trên các bộ dữ liệu khác như RAVDESS, EmoDB, MSP-IMPROV, đồng thời tiến hành xây dựng và thu thập bộ dữ liệu cảm xúc giọng nói chuẩn dành cho tiếng Việt. Định hướng thứ hai tập trung vào cải tiến cấu trúc mô hình bằng cách tích hợp các cơ chế chú ý (Self-Attention, Multi-Head Attention), nghiên cứu áp dụng các kiến trúc Transformer cho bài toán SER và khai thác các mô hình pre-trained quy mô lớn như Wav2Vec 2.0 hay HuBERT cho bước trích xuất đặc trưng. Định hướng thứ ba liên quan đến phát triển hệ thống đa phương thức (multimodal) kết hợp thông tin âm thanh, văn bản và hình ảnh khuôn mặt thông qua các cơ chế Cross-Modal Attention để tối ưu hóa độ chính xác. Định hướng thứ tư là tối ưu hóa mô hình MobileNet-V2 gọn nhẹ nhằm triển khai ứng dụng thời gian thực trên các thiết bị di động và tích hợp vào hệ thống tổng đài thông minh Call Center. Định hướng thứ năm là thực hiện nghiên cứu so sánh đối chứng bằng việc huấn luyện lại mô hình theo phương pháp phân chia phụ thuộc người nói (Speaker-Dependent split), giúp đánh giá chi tiết mức độ suy giảm hiệu năng thực tế và làm nổi bật những thách thức cốt lõi khi chuyển đổi từ cấu hình phụ thuộc người nói sang cấu hình độc lập người nói (Speaker-Independent) đã được triển khai trong đồ án này. Định hướng cuối cùng là cải tiến quy trình huấn luyện thông qua Curriculum Learning, tự giám sát (Self-supervised learning) và thích ứng miền (Domain Adaptation) để chuyển giao hiệu quả tri thức giữa các ngôn ngữ khác nhau.
